\chapter{Grade–tonnage}\label{ch:gt}
\begin{goals}
\textbf{In\index{grade–tonnage} this chapter} you draw tonnage and mean grade above cutoff for a block model, give it a density and units so the tonnes are tonnes, overlay a second model or a second domain, read a reconciliation as two curves, and write the same curve from a routine. On a drillhole layer the same window adds up interval length in meters instead.
\end{goals}

Grade–tonnage answers the first question anyone asks of a block model: \emph{how much is there above a given grade, and at what grade?} micro computes the curve exactly, at any model size, from the source file — every block weighted by its own volume and density.

Open it from a block model's (or a drillhole layer's) right-click menu, \menu{Analyse>Grade-tonnage…}, or press \key{⌘K} (\key{Ctrl}+\key{K} on Windows and Linux) and type \code{grade}. The window works as \cref{ch:analysis} describes; this chapter covers what is particular to grade–tonnage.

\section{A first curve}\label{sec:gt-first}
The demo project opens with a grade–tonnage window already run on \file{deposit.csv}, floating at the right of the view. It is set up the way a report needs it: \col{FE} as the grade, declared in \%, and \col{SG} as the density. So its left axis reads \emph{tonnage (t)}, and the subtitle \emph{tonnes · 40 cutoffs · 1,786,620 rows}.

\begin{tryit}
Open the demo project and find the grade–tonnage window at the right of the view. (If you closed it, right-click \file{deposit.csv}, choose \menu{Analyse>Grade-tonnage…} and press \ui{Run}.)
\begin{enumerate}
\item Read the series box: the layer \file{deposit.csv}, the grade column \col{FE} with its unit \ui{\%} beside it, and \col{SG} in the field below — \emph{density expr (blank = count)}. Under \ui{display}, the plot \ui{unit} is \ui{\%} too.
\item Change \ui{cutoffs} from 40 to 20. The curve redraws at once — no new pass.
\item Clear the density field and press \key{Enter}. The window notes \emph{↻ config changed — Run to update}; the old curve stays until you act.
\item Press \ui{Run}. The left axis now reads \emph{volume (m³)} and the subtitle \emph{volume (no density)}: the same curve in cubic meters.
\item Type \code{SG} back into the density field and press \ui{Run} again.
\item Click \ui{⧉ table} and paste into a spreadsheet: one row per cutoff, with tonnes, grade and contained metal in tonnes.
\end{enumerate}
\end{tryit}

\shot{gt-window}{Grade–tonnage on the demo deposit with \col{FE} as the grade and \col{SG} as the density: tonnage above cutoff (solid, left axis) falls while the mean grade of what remains (dashed, right axis) rises.}

The solid line is tonnage on the left axis; the dashed line of the same color is the mean grade of that tonnage, on the right axis. The demo deposit is mostly waste around an iron-ore lens, and the curve shows the two populations: the tonnage falls steeply between 20 and 30\,\% — 2.28~Gt grades above 20\,\%, only 0.32~Gt above 30\,\% — and then declines slowly through the ore. \Cref{fig:gt-demo} is the ore part of the curve, at six stated cutoffs, drawn from micro's numbers: 322~Mt at 47.9\,\% above 30, 173~Mt at 55.4\,\% above 50, 24~Mt at 62.3\,\% above 60.

\begin{figure}[htbp]\centering
\begin{tikzpicture}
\begin{axis}[width=0.86\linewidth, height=58mm, xlabel={cutoff (\% Fe)}, ylabel={tonnes above cutoff (Mt)}, ymin=0, legend style={at={(0.03,0.05)}, anchor=south west}]
\addplot[amber, very thick, mark=*, mark size=1.4pt] table[col sep=tab, x=cutoff, y expr=\thisrow{tonnes}/1e6]{figures/data/gt-demo.tsv};
\addlegendentry{tonnage}
\addlegendimage{info, very thick, dashed, mark=square*, mark size=1.3pt}
\addlegendentry{mean grade}
\end{axis}
\begin{axis}[width=0.86\linewidth, height=58mm, axis y line*=right, axis x line=none, grid=none, ylabel={mean grade above cutoff (\% Fe)}]
\addplot[info, very thick, dashed, mark=square*, mark size=1.3pt] table[col sep=tab, x=cutoff, y=grade]{figures/data/gt-demo.tsv};
\end{axis}
\end{tikzpicture}
\caption{The demo deposit's grade–tonnage curve for \col{FE}, tonnage from \col{SG} × block volume, at the cutoffs 30, 40, 50, 55, 60 and 65\%.}\label{fig:gt-demo}
\end{figure}

\section{How the numbers are made}\label{sec:gt-method}
For each series micro makes one pass over the layer. For every record it evaluates the grade $g$, the density $\rho$ and the block's volume $V$, and forms the weight
\[ w = \rho \times V . \]
Then, for a cutoff $c$, tonnage is the sum of $w$ over the records with $g \ge c$, the mean grade is $\sum w g / \sum w$ over the same records, and contained metal is tonnage × mean grade.

Four details decide what goes into those sums. The first three are the shared rules of \cref{sec:analysis-rules}, repeated here because in this window they decide the tonnage you report.

\paragraph{The volume is the block's own.} For a regular model $V$ is the lattice cell: 10 × 10 × 5 = 500 m³ in the demo. For a sub-blocked model it is each record's own size, read from its dimension columns, so a parent block and a sub-block weigh what they are (\cref{ch:blocks}). Weighting every record by the fine cell would count a 20~m parent as one 2.5~m sub-block and lose most of the waste tonnage.

\paragraph{A blank grade is unestimated.} A record whose grade is blank or not numeric is skipped: it adds no tonnage at any cutoff, including the lowest. That is the right reading of a block the estimator never reached. It also means the tonnage at the lowest cutoff is the tonnage of the \emph{estimated} model, not of the whole lattice. If you need the unestimated tonnage as a line of its own, the cross-tab keeps it (\cref{ch:resources}).

\paragraph{A zero weight is no weight.} A record whose density evaluates to blank, zero or a negative number is skipped too. Blank density is the commonest reason a tonnage comes out low.

\paragraph{“Above” includes the cutoff.} A block grading exactly 55 is counted at a cutoff of 55.

\begin{why}
micro does not sort the model, and it does not keep every grade in memory. The pass drops each record into a fine grade step — a power of ten chosen so that at least 20,000 steps span the data, 0.001 for the demo's \col{FE} — and keeps the exact sum of $w$ and of $w g$ per step. Tonnage and grade at a cutoff that lies \emph{on} a step are then exact, with memory that depends on the number of occupied steps, not on the number of blocks. That is why a 60-million-block model gives the same curve as a spreadsheet would, if a spreadsheet could hold it.
\end{why}

For the cutoffs you state in a recipe or routine (\cref{sec:gt-routine}), the consequence is simple: any cutoff you would normally type — 0.5, 55, 0.35 — lies on a step and is exact. A cutoff \emph{between} two steps, such as 40.0001 on a 0.001 step, is counted from the next step up, and micro says so in a warning (\emph{… not on the 0.001 grade step — counted from the next step up}).

Exact arithmetic is not a correct model. The curve describes the blocks as estimated, at their support, and it inherits the estimate's smoothing: an over-smoothed model shows too little grade above a high cutoff, and nothing in this window corrects that. Testing the model against its data is the subject of \cref{ch:validation}.

\section{The settings}\label{sec:gt-settings}
\begin{tabularx}{\linewidth}{@{}lMD@{}}\toprule
setting & meaning & default \\ \midrule
series · layer & the layer this curve reads & the window's layer \\
series · column & the grade column; calculated and stored columns are offered too & the first numeric column \\
series · unit & what the grade is in: \%, g/t, ppm, oz/t & the layer's declaration \\
series · density & an expression: \code{SG}, \code{2.7}, \code{SG * PARTIAL}. Blank = volume & the layer's, else blank \\
series · filter & rows for this series only, e.g. \code{LITO = "HEMATITE"} & none \\
\ui{cutoffs} & how many cutoff steps to draw, 5–200 & 40 \\
\ui{unit} & the plot unit: the cutoff and mean-grade axes & the layer's declaration \\
\ui{ρ unit} & the unit of the density expressions: t/m³ (= g/cm³), kg/m³, lb/ft³ & t/m³ \\
\ui{scope} & \ui{all} · \ui{filt} · \ui{sel} (\cref{sec:analysis-scope}) & all \\
ranges & manual \ui{min}/\ui{max} for the cutoff, tonnage and grade axes & automatic \\ \bottomrule
\end{tabularx}
\medskip

Changing a series' layer, column, density or filter, or the \ui{ρ unit}, needs a new \ui{Run}. The cutoff count, the units of grade and the ranges are live.

\section{Density}\label{sec:gt-units}
The\index{density} density field is an expression in the same language as the filter (\cref{app:expr}), evaluated per block:

\begin{itemize}
\item \code{SG} — a density column;
\item \code{2.7} — a constant, for a model without one;
\item \code{SG * PARTIAL} — a density scaled by an ore fraction or a partial-block percentage. A fraction column made by a solid query (\cref{ch:geometry}) works here, which gives the tonnage inside a wireframe;
\item \code{if(LITO = "CANGA", 3.1, SG)} — a correction by rock type.
\end{itemize}

Leave it blank and the weight is the block volume alone: the axis is labeled \emph{volume (m³)} and the subtitle says \emph{volume (no density)}. Those two labels are the only warnings in this window, so read them before you read a number off the curve.

\begin{caution}
A density of 2,700 with the \ui{ρ unit} left at t/m³ gives tonnages a thousand times too large, and nothing can detect it — 2,700 t/m³ is a number like any other. If your density column is in kg/m³, declare it in \ui{ρ unit}. Model QA (\cref{ch:qa}) flags densities outside 1–6 t/m³.
\end{caution}

Grade units work as in every analysis window (\cref{sec:analysis-units}): declared on the series and on the plot, they put series in different units on one honest axis and turn the table's last column into tonnes of metal — or m³ of metal, with no density.

\section{Cutoffs}\label{sec:gt-cutoffs}
In\index{cutoff} the window you choose \emph{how many} cutoff steps, not which cutoffs. micro draws one more round value than the count — 40 steps are 41 cutoffs, the first and the last included — evenly spaced by a step of 1, 1.2, 1.5, 2, 2.5, 3, 4, 5, 6 or 8 × a power of ten, from a round number at or below the lowest grade to one at or above the highest. Every one is on a grade step, so every plotted point is exact.

Automatic cutoffs are good for looking and wrong for reporting: they come from the data's range, so adding one high-grade block can move them all. A routine therefore requires an explicit list (\cref{sec:gt-routine}), and a figure meant to be compared with last quarter's should come from a recipe or routine with stated cutoffs, or at least fixed axis ranges.

To read a single cutoff precisely in the window, copy the table (\ui{⧉ table}) — it has one row per cutoff drawn — or raise the cutoff count until your value is among the steps.

\section{Comparing models}\label{sec:gt-compare}
Each\index{grade–tonnage!comparing models} series picks its own layer, so you compare two models, or two domains of one, in a single window.

\begin{tryit}
With the demo's grade–tonnage window open (right-click \file{deposit.csv}, \menu{Analyse>Grade-tonnage…}, if you closed it):
\begin{enumerate}
\item Click \ui{+ add series}. A second series appears on the window's layer, with its first numeric column (\col{FE}) and an empty filter.
\item In its density field type \code{SG}.
\item In its filter field type \code{LITO = "HEMATITE"}.
\item Press \ui{Run}. Two pairs of curves appear: the deposit, and its hematite — flat at about 24~Mt until the cutoff passes 60\,\%, then running along the deposit's curve.
\item Click a series' colored dot to change its color. Both series read the same layer and column, so the legend names both \emph{FE}; the color tells them apart.
\end{enumerate}
\end{tryit}

With two block models loaded — this estimate and the previous one, kriging and nearest-neighbor — set the second series' layer menu to the other model and choose its column. The legend then names each curve by layer (without its extension) and column: \emph{model\_A·FE · t} and \emph{model\_A·FE · grade}, then the same for \emph{model\_B}.

The two models need not share a lattice, a block size or even a format, because each series is weighed by its own layer's block volumes. So when a re-blocked model changes the curve, you are looking at the support effect — larger blocks, less variance, fewer tonnes above a high cutoff — not at an artifact of the tool.

\section{On a drillhole layer}\label{sec:gt-holes}
Grade–tonnage also runs on a drillhole interval table: right-click the drillhole layer and choose \menu{Analyse>Grade-tonnage…}, or make it active, press \key{⌘K} and type \code{grade}. The window is the same; what it adds up is not.

A drillhole interval has no block volume. It weighs its length, TO − FROM, so the curve is \emph{meters of drilling above the cutoff}, on an axis labeled \emph{length (m)}, with the subtitle \emph{length-weighted (m)}. The mean grade is length-weighted, as a composite's would be, and the table's last column is the \term{accumulation}, grade × length (\%·m, g/t·m) — the number a drill intercept is quoted in.

A density multiplies the length: the axis becomes \emph{length × density (m·t/m³)} and the subtitle says \emph{not tonnes — no section area}. That product is tonnes per square meter of section; an interval table holds no area, so micro never calls it tonnes. On the demo's \file{demo\_assay.csv}, 2,205~m of the drilling grades 30\,\% Fe or better, at a length-weighted 48.7\,\%; 80~m grades 60\,\% or better.

\begin{caution}
A block model and a drillhole layer cannot share one curve: cubic meters and meters of core are not the same measure. A block model's window lists only block models in its series' layer menu, a drillhole layer's window only drillhole layers, and a routine step that mixes them stops with \emph{series measure different things — block models accumulate volume / tonnes (…), drillholes accumulate interval length in metres (…)}. Run them as two curves. For a model-against-samples check use Q-Q (\cref{ch:stats}) or the validation windows (\cref{ch:validation}), which are made for it.
\end{caution}

\section{Reconciliation-aware curves}\label{sec:gt-reconcile}
When\index{reconcile!grade–tonnage} you reconcile two models (\cref{ch:joins}), micro makes a derived layer carrying both grades side by side on the common lattice, in columns named \col{ref\_}\emph{grade} and \col{this\_}\emph{grade} — \col{ref\_FE} and \col{this\_FE}.

Open grade–tonnage on that layer and it recognizes the pair: the window starts with two series, the reference and the new model, instead of one. Press \ui{Run} and you have the reconciliation as a pair of grade–tonnage curves over exactly the same blocks. The Δ map of \cref{ch:joins} shows \emph{where} the models differ; this shows what the difference does to the tonnage above each cutoff.

Because both series read one layer, the scope applies to both. Select the bench that was mined this month, set the scope to \ui{sel}, and the pair of curves is the reconciliation of that bench alone, recomputed whenever the selection changes.

\section{Before you quote a number}\label{sec:gt-pitfalls}
\begin{itemize}
\item The left axis says \emph{tonnage (t)}, not \emph{volume (m³)} — no density, no tonnes (\cref{sec:gt-units}).
\item The density is in the declared \ui{ρ unit}, and it has no gaps: a blank density drops the block (\cref{sec:gt-method}).
\item Unestimated blocks are in no cutoff; the cross-tab's \ui{no grade} column says how much that is (\cref{sec:xtab-nograde}).
\item There is no \emph{↻ config changed} note under \ui{Run} (\cref{sec:analysis-free}).
\item The subtitle does not start with \emph{selection} or \emph{filtered} unless you meant it (\cref{sec:analysis-scope}).
\item The cutoffs in a report are stated, not automatic (\cref{sec:gt-cutoffs}).
\end{itemize}

\section{In a routine}\label{sec:gt-routine}
The tool is \code{gt} (its settings: \cref{tool:gt}). They are the window's, with one addition — \code{cutoffs}, the explicit list, in the plot unit \code{gradeUnit} — which a routine must give.

\begin{lstlisting}[language=yaml]
- tool: "gt"
  with:
    layer: "deposit.csv"
    gradeUnit: "pct"
    densUnit: "t/m3"
    series:
      - col: "FE"
        unit: "pct"
        weight: "SG"
    cutoffs:
      - 30
      - 40
      - 50
      - 55
      - 60
      - 65
  as: "gt"
- check:
    every: "tonnes >= 0"
    in: "gt"
  say: "no negative tonnage"
\end{lstlisting}

The report is a machine table with one row per cutoff and series, keyed by \col{series} and \col{cutoff}:

\medskip
\begin{tabularx}{\linewidth}{@{}lL@{}}\toprule
column & meaning \\ \midrule
\col{series} & the grade column; with several series, \emph{layer·column}, the layer without its extension (\code{deposit·FE}) \\
\col{cutoff} & the cutoff, in \code{gradeUnit} \\
\col{tonnes} & tonnage at or above the cutoff — named \col{volume} (m³) when no series has a density; on a drillhole layer \col{length} (m), or \col{weighted\_length} (m·t/m³) with a density \\
\col{grade} & the mean grade of that tonnage; blank when there is none \\
\col{metal} & tonnes × grade, in tonnes of metal when \code{gradeUnit} is declared (m³ of metal without a density); on a drillhole layer \col{accumulation}, grade × length \\ \bottomrule
\end{tabularx}
\medskip

\begin{caution}
Two series that read the same column of the same layer get the same name in the \col{series} column, so their rows cannot be told apart in the report. In a routine, give each such comparison its own step — one \code{gt} step per filter, each with its own \code{as:} name.
\end{caution}

A block model opened as points (\cref{ch:opening}) has no block volumes to weigh, so a \code{gt} step on it stops with \emph{deposit.csv: no block grid — a model opened as points has no block volumes to weigh by (reopen it as a block model)}, rather than counting every point as a cubic meter.

A cutoff that is not on the grade step is not an error: the step runs, and the warning appears on the step's line in the run tracker. The resolved settings in the run record (\cref{ch:records}) list the cutoffs actually used, sorted ascending.

To judge the curve, add checks (\cref{ch:checks}): a row check as above, or a comparison of the whole table with the accepted baseline run — \code{compare: "gt"} with a \code{tolerance} — which is the month-over-month test that a model's tonnage above each cutoff has not moved by more than you allow.
